\documentclass[11pt]{article}
\usepackage[T1]{fontenc}
\usepackage{lmodern,amsmath,amssymb,amsthm,mathtools,microtype}
\usepackage[a4paper,margin=25mm]{geometry}
\usepackage[hidelinks]{hyperref}
\setlength{\emergencystretch}{2em}
\newtheorem{theorem}{Theorem}[section]
\newtheorem{lemma}[theorem]{Lemma}
\newtheorem{corollary}[theorem]{Corollary}
\theoremstyle{remark}\newtheorem{remark}[theorem]{Remark}
\newcommand{\Ric}{\operatorname{Ric}}
\newcommand{\tr}{\operatorname{tr}}
\newcommand{\Id}{\operatorname{Id}}
\newcommand{\RR}{\mathcal R_+}
\title{A Determinant-Sign Case of the Four-Dimensional\\Negative Goldberg Problem}
\author{}
\date{13 October 2026}
\begin{document}
\maketitle
\begin{abstract}
A closed negative Einstein almost-K\"ahler four-manifold is K\"ahler if
$\det(W^++s\Id/12)\le0$ everywhere in its almost-complex orientation.
The proof combines the classical Einstein Weyl--Weitzenb\"ock equation with
Armstrong's existence of a torsion zero. A scalar maximum principle first
forces the shifted curvature operator to be nonpositive. The torsion zero
then places a second elementary symmetric function on the boundary of its
positive cone; the strong maximum principle forces a parallel rank-one
operator and finally the parallelness of the given fundamental form.
The determinant assumption is additional: the general negative Goldberg
problem remains the case not covered by this sign condition.
\end{abstract}

\section{Statement and normalization}
All manifolds below are smooth, closed (compact without boundary), connected,
and four-dimensional. Curvature on two-forms is normalized by
\[
 \mathcal R=\begin{pmatrix}W^++s\Id/12&B^*\\ B&W^-+s\Id/12\end{pmatrix}.
\]
Weyl norms are Hilbert--Schmidt norms of endomorphisms. The scalar and rough
Laplacians are nonnegative: $\Delta=-\operatorname{div}\nabla$ and
$\nabla^*\nabla$. Set
\[
 \Ric=-a g,\qquad a>0,\qquad s=-4a,\qquad
 \RR=W^+-\frac a3\Id.
\]

\begin{theorem}\label{thm:main}
Let $(M,g,J,\Omega)$ be almost-K\"ahler and Einstein with $\Ric=-a g$,
$a>0$. Choose the orientation induced by $J$. If
\[
                 \det\RR\le0\quad\hbox{everywhere on }M,
\]
then $\nabla\Omega=0$. In fact, for $n=\Omega/\sqrt2$,
\[
                    \RR=-a\,n\otimes n.
\]
\end{theorem}

The hypothesis involves the full self-dual curvature block, not merely
$W^+$, not the determinant of the full six-dimensional curvature operator,
and not sectional-curvature negativity. The proof below uses neither
project FD/N1 estimates nor a uniform compactness assumption over a family
of metrics. The inputs are classical identities and Armstrong's topological
observation; no literature priority is asserted for their composition here.

\section{The scalar equation that supplies the sign}
Put
\[
 A=-\RR=\frac a3\Id-W^+,\qquad \tr A=a,
 \qquad u=\sigma_2(A)=\frac{a^2-|A|^2}{2}
                  =\frac{a^2}{3}-\frac12|W^+|^2.
\]
These definitions are smooth even at repeated curvature eigenvalues.

\begin{lemma}\label{lem:scalar}
For every closed oriented Einstein four-manifold with $\Ric=-ag$, one has
\begin{equation}\label{eq:scalar}
 \boxed{(\Delta+2a)u=|\nabla A|^2+18\det A.}
\end{equation}
\end{lemma}
\begin{proof}
Write $W=W^+$. The classical Einstein Weyl--Weitzenb\"ock equation is
\begin{equation}\label{eq:weyl}
 \nabla^*\nabla W+\frac s2W-6W^2+2|W|^2\Id=0.
\end{equation}
See Derdzi\'nski~\cite{Derdzinski}; the same operator convention is written
explicitly in \cite[Lemma 2.5, $f=1$]{CD}.
Since $\tr W=0$ and $\tr(W^3)=3\det W$,
\[
 \Delta|W|^2=-s|W|^2+36\det W-2|\nabla W|^2.
\]
The shift $W=a\Id/3-A$ gives
\[
 |W|^2=\frac{2a^2}{3}-2u,
 \qquad
 \det W=-\frac{2a^3}{27}+\frac a3u-\det A.
\]
As $\Delta u=-\tfrac12\Delta|W|^2$, substitution gives
$\Delta u=-2au+18\det A+|\nabla A|^2$, as asserted.
\end{proof}

\begin{lemma}\label{lem:cone}
If $\det\RR\le0$ everywhere, then $A\succeq0$ everywhere.
If, in addition, $u(p)=0$ at any point, then
\[
 u\equiv0,\qquad \nabla A=0,\qquad A^2=aA,
 \qquad\operatorname{rank}A=1.
\]
\end{lemma}
\begin{proof}
Since the block has odd rank, the determinant assumption is precisely
$\det A\ge0$. Equation~\eqref{eq:scalar} implies
$(\Delta+2a)u\ge0$. A negative minimum of $u$ would satisfy
$\Delta u\le0$ and $2au<0$, a contradiction. Hence $u\ge0$.

We use the following elementary fact: a real symmetric $3\times3$ matrix
with positive trace, nonnegative $\sigma_2$, and nonnegative determinant
is positive semidefinite. If there is exactly one negative eigenvalue and
the other two are positive, the determinant is negative. If they are
$-x,-y,a+x+y$ with $x,y>0$, then
\[
 \sigma_2=xy-(x+y)(a+x+y)<0.
\]
If one eigenvalue is zero, a negative eigenvalue and positive trace likewise
force $\sigma_2<0$. The remaining cases are incompatible with positive
trace. Applying this fact proves $A\succeq0$.

If $u(p)=0$, the strong minimum principle, applied to the nonnegative
function $u$ with $(\Delta+2a)u\ge0$, gives $u\equiv0$.
Equation~\eqref{eq:scalar} then gives $\nabla A=0$ and $\det A=0$.
A positive semidefinite matrix of trace $a$ and $\sigma_2=0$ has eigenvalues
$(a,0,0)$, proving the remaining assertions.
\end{proof}

Thus, on this closed negative Einstein class, the global determinant-sign
condition is equivalent to $\RR\preceq0$. This equivalence uses the
Einstein equation and the maximum principle; it is false for arbitrary
pointwise symmetric matrices with negative trace.

\section{A torsion zero forces the rigid boundary case}
We recall the classical global input in the precise form needed.

\begin{lemma}[Armstrong]\label{lem:zero}
A closed Einstein almost-K\"ahler four-manifold with $s\ne0$ has a point
where $\nabla\Omega=0$.
\end{lemma}
\begin{proof}
The Nijenhuis tensor is a section of the complex rank-two bundle
\[
 \mathcal N=\Lambda^{0,2}T^*M\otimes T^{1,0}M.
\]
Its second Chern number is
\[
 c_2(\mathcal N)[M]=2c_1(TM,J)^2+c_2(TM,J)=5\chi+6\tau.
\]
If it were nowhere zero, this number would vanish. On an Einstein
four-manifold with nonzero scalar curvature, Chern--Weil gives
\[
 8\pi^2\chi=\int\left(|W^+|^2+|W^-|^2+\frac{s^2}{24}\right)>0,
\]
\[
 4\pi^2(2\chi+3\tau)=\int\left(2|W^+|^2+\frac{s^2}{24}\right)>0.
\]
Consequently $5\chi+6\tau=\chi+2(2\chi+3\tau)>0$, a contradiction.
On an almost-K\"ahler manifold, vanishing of the Nijenhuis tensor at a point
is equivalent to vanishing of $\nabla\Omega$ there. This is
\cite[Theorem 2.2.4]{Armstrong}, whose topological proof is reproduced here.
\end{proof}

\begin{proof}[Proof of Theorem~\ref{thm:main}]
Let $n=\Omega/\sqrt2$. It is a harmonic self-dual two-form of unit length.
The harmonic-form Weitzenb\"ock equation gives
\begin{equation}\label{eq:n}
 \nabla^*\nabla n=2W^+n-\frac s3n=2(a\Id-A)n.
\end{equation}
Since $|n|=1$, putting $q=|\nabla\Omega|^2=2|\nabla n|^2$ yields
\begin{equation}\label{eq:rayleigh}
                     q=4\bigl(a-\langle An,n\rangle\bigr).
\end{equation}
Choose $p$ from Lemma~\ref{lem:zero}. Lemma~\ref{lem:cone} gives $A\succeq0$;
equation~\eqref{eq:rayleigh} gives
$\langle A(p)n(p),n(p)\rangle=a=\tr A(p)$.
It follows that
\[
 A(p)=a\,n(p)\otimes n(p),\qquad u(p)=0.
\]
The second part of Lemma~\ref{lem:cone} now applies. In particular,
$\Pi=A/a$ is a globally defined parallel orthogonal rank-one projection.

Define $v=\Pi n$ without making any orientation choice for the image line
bundle. By~\eqref{eq:n} and $A^2=aA$,
\[
 \nabla^*\nabla v=\Pi\nabla^*\nabla n
                =2\Pi(a\Id-A)n=0.
\]
Integration over the closed manifold gives $\nabla v=0$.
At $p$, $v(p)=n(p)$, so $|v|\equiv1$. Since $|n|\equiv1$ and $\Pi$ is an
orthogonal projection, $n=v$ everywhere. Therefore $\nabla n=0$, hence
$\nabla\Omega=0$. The relation $\RR=-a\,n\otimes n$ follows as well.
\end{proof}

\section{Actual consequences and the remaining case}
\begin{corollary}\label{cor:positive}
Every closed strictly almost-K\"ahler negative Einstein four-manifold would
have an open set on which $\det\RR>0$. On that set the signature of
$\RR$ is $(+,-,-)$.
\end{corollary}
\begin{proof}
The first statement is the contrapositive of the theorem and continuity.
The determinant is positive, while $\tr\RR=-a<0$, so exactly one eigenvalue
is positive and the other two are negative.
\end{proof}

\begin{corollary}\label{cor:definite}
A closed negative Einstein four-manifold with $\RR\prec0$ everywhere has
no compatible almost-K\"ahler structure in the stated orientation.
\end{corollary}
\begin{proof}
Such a structure would imply $\RR=-a\,n\otimes n$, which has two zero
eigenvalues. Alternatively, strict negativity makes the right side of
\eqref{eq:rayleigh} positive everywhere, contradicting Lemma~\ref{lem:zero}.
\end{proof}

Fine--Premoselli construct compact, negatively curved, non-locally-homogeneous
Einstein four-metrics and explicitly state that their examples satisfy
$W^++s\Id/12\prec0$~\cite[Theorem 1.2 and following discussion]{FP}.
Corollary~\ref{cor:definite} applies to those metrics in that orientation.
This concerns compatibility with the specified metrics, not nonexistence
of symplectic forms or Einstein metrics on the underlying smooth manifolds.
No claim about the opposite orientation is needed here.

In the preceding harmonic-product representation, this says that under the
determinant condition every positive trace-one rank-one element of the
harmonic multiplication kernel comes from a parallel form. In the strictly
negative-definite case there is no such rank-one element at all. The
Einstein PDE, rather than positivity of the finite-dimensional matrices
alone, supplies this conclusion.

\paragraph{Scope.}
The original Q4 hypotheses do not include $\det\RR\le0$. No derivation of
that sign from those hypotheses is given here. The proof settles this
specified curvature-sign case and identifies a necessary open
indefinite-curvature region in any hypothetical strict closed example.
It neither proves full negative Goldberg nor supplies a closed counterexample.

\begin{thebibliography}{9}
\bibitem{Derdzinski}
A. Derdzi\'nski,
\emph{Self-dual K\"ahler manifolds and Einstein manifolds of dimension four},
Compositio Math. \textbf{49} (1983), 405--433.
\url{https://archive.numdam.org/item/CM_1983__49_3_405_0/}.
\bibitem{CD}
G. Catino and D. Dameno,
\emph{Conformal K\"ahler Rigidity of Einstein Four-Manifolds},
arXiv:2607.21538v1 (2026), Lemma 2.5.
\url{https://arxiv.org/html/2607.21538v1}.
Only the classical Weyl--Weitzenb\"ock identity is used.
\bibitem{Armstrong}
J. Armstrong, \emph{Almost-K\"ahler Geometry}, D.Phil. thesis, Oxford (1998),
Theorem 2.2.4.
\url{https://johnarmstrongmaths.com/publications/resources/thesis.pdf}.
\bibitem{FP}
J. Fine and B. Premoselli,
\emph{Examples of compact Einstein four-manifolds with negative curvature},
arXiv:1802.00608, Theorem 1.2 and the curvature inequality following it.
\url{https://arxiv.org/html/1802.00608}.
\end{thebibliography}
\end{document}
